\chapter{Neural Networks}
\label{ch:neuralnets}

\begin{tcolorbox}[feynbox={Sky}{BBg},title={Before and After}]
\textbf{Before this chapter you need:} composition of functions
(\S\ref{sec:composition}), the chain rule
(\S\ref{sec:rules}, \S\ref{sec:chainmulti}), gradients and Jacobians
(\S\ref{sec:gradient}, \S\ref{sec:jacobian}), matrices and matrix
multiplication (\S\ref{sec:vectors}--\S\ref{sec:matmul}), gradient descent and
batching (\S\ref{sec:gradientdescent}, \S\ref{sec:sgd}).\\
\textbf{After it you will understand:} what a neuron computes, why the choice
between sigmoid and $\tanh$ changed this project's results by more than any other
single edit, how backpropagation is the chain rule run backwards, and the two
architectures the models here are built from --- gated recurrent cells and dilated
causal convolutions.
\end{tcolorbox}

A neural network is a composition of simple functions with adjustable
coefficients. Nothing in this chapter is conceptually new; it is
Chapter~\ref{ch:calculus2}'s chain rule and Chapter~\ref{ch:linalg}'s matrices,
assembled and given names.

%% ─────────────────────────────────────────────────────────────────────────────
\section{A Neuron}
\label{sec:neuron}

\situation{
A loan officer weighs several facts --- income, debts, employment history --- giving
each a different importance, adds them up, and issues a yes or no. The weighting
is the judgement; the threshold turns a score into a decision.
}

\problem{
We need the smallest unit that can be composed into something expressive and
adjusted by gradient descent. It must combine many inputs into one number, and it
must be differentiable so \S\ref{sec:gradientdescent} applies.
}

\workedarith{
Inputs $\mathbf{x} = [2, 3, 1]$, weights $\mathbf{w} = [0.5, -1, 2]$, bias
$b = 0.4$.

\medskip
\textbf{Weighted sum.} This is the dot product of \S\ref{sec:vectors}:
\[
  z = \mathbf{w}\cdot\mathbf{x} + b = 0.5(2) + (-1)(3) + 2(1) + 0.4
    = 1 - 3 + 2 + 0.4 = 0.4 .
\]

\textbf{Activation.} Apply $\tanh$:
\[
  a = \tanh(0.4) = 0.3799 .
\]

\medskip
\textbf{What the weights encode.} The second input has weight $-1$: raising it
lowers the output. The third has weight $2$: it matters twice as much per unit as
the first. Learning means adjusting these numbers.

\medskip
\textbf{Why the bias is needed.} Without $b$, an all-zero input forces $z = 0$ and
the neuron cannot represent a non-zero default. The bias is the intercept $c$ of
\S\ref{sec:linear}, and it is what makes the map affine rather than strictly
linear (\S\ref{sec:vectorspaces}).
}

\formaldef{
A \textbf{neuron} computes
\[
  a = \sigma\!\left(\mathbf{w}\cdot\mathbf{x} + b\right)
    = \sigma\!\left(\sum_{i} w_i x_i + b\right) ,
\]
where $\mathbf{w}$ are \textbf{weights}, $b$ is the \textbf{bias}, and $\sigma$ is
a non-linear \textbf{activation function}.

A \textbf{layer} applies many neurons to the same input, which is one matrix
multiplication:
\[
  \mathbf{a} = \sigma(W\mathbf{x} + \mathbf{b}) ,
\]
with $\sigma$ applied entrywise. For $n$ inputs and $m$ neurons, $W$ is
$m \times n$.
}

\analogybreak{
The biological analogy is thin and worth dropping early. Real neurons fire
discrete spikes, have complex temporal dynamics, and do not learn by
backpropagation --- there is no known mechanism for propagating an error signal
backwards through a brain. The name is historical. What is being described is a
parameterised differentiable function, and nothing about its usefulness depends on
the biology.
}

\whyhere{
Every parameter counted in this project is a weight or a bias in such a unit. The
TimeGAN embedder's $9{,}744$ parameters and its generator's $11{,}400$ are counts
of exactly these numbers, and gradient descent adjusts every one of them at each
step.
}

%% ─────────────────────────────────────────────────────────────────────────────
\section{Activation Functions and Their Derivatives}
\label{sec:activations}

\situation{
A stack of currency conversions is still a single conversion
(\S\ref{sec:matmul}). Stack fifty and you have one exchange rate. Something has to
break the linearity, or depth buys nothing at all.
}

\problem{
The non-linearity must do two jobs at once: make composition expressive, and pass
a usable gradient backwards. These pull against each other, and getting the
balance wrong silently cripples training --- as it did here.
}

\workedarith{
\textbf{The three functions and their derivatives.}
\[
  \sigma(x) = \frac{1}{1+e^{-x}} , \qquad \sigma'(x) = \sigma(x)\bigl(1 - \sigma(x)\bigr) ,
\]
\[
  \tanh(x) = \frac{e^{x}-e^{-x}}{e^{x}+e^{-x}} , \qquad \tanh'(x) = 1 - \tanh^2(x) ,
\]
\[
  \mathrm{ReLU}(x) = \max(0,x) , \qquad \mathrm{ReLU}'(x) = \begin{cases} 1 & x>0 \\ 0 & x<0 \end{cases}
\]

\medskip
\textbf{At the centre of each range.} Sigmoid outputs $0.5$ when $x = 0$, and
there
\[
  \sigma'= 0.5 \times 0.5 = 0.25 .
\]
$\tanh$ outputs $0$ when $x = 0$, and there
\[
  \tanh' = 1 - 0^2 = 1.00 .
\]

\medskip
\textbf{A factor of four, per layer.} Backpropagation multiplies these along the
chain (\S\ref{sec:chainmulti}). Through three layers, at the most favourable
point of each:

\medskip
\begin{tabular}{@{}lrrr@{}}
\toprule
Layers & sigmoid & $\tanh$ & ratio \\
\midrule
1 & $0.25$    & $1.0$ & $4$ \\
2 & $0.0625$  & $1.0$ & $16$ \\
3 & $0.0156$  & $1.0$ & $64$ \\
\bottomrule
\end{tabular}

\medskip
Through a three-layer stack the sigmoid passes at most $1.6\%$ of the gradient
signal that $\tanh$ passes --- and $0.25$ is its \emph{best} case. Away from the
centre it is far worse: at $x = 3$, $\sigma(3) = 0.953$ and
$\sigma' = 0.953(0.047) = 0.045$.

\medskip
\textbf{Saturation.} Both sigmoid and $\tanh$ flatten at their extremes, where the
derivative approaches zero and learning stalls. ReLU has derivative exactly $1$
for all positive inputs and so does not attenuate at all --- at the cost of
returning exactly $0$ for negative inputs, where the unit can stop learning
permanently (a ``dead'' unit).
}

\formaldef{
An \textbf{activation function} is a non-linear scalar function applied entrywise.
Ranges and uses:
\begin{itemize}
  \item $\sigma$: range $(0,1)$, for probabilities and gates;
  \item $\tanh$: range $(-1,1)$, zero-centred, for hidden and bounded outputs;
  \item ReLU: range $[0,\infty)$, cheap and non-saturating, the default in deep
        feedforward and convolutional stacks.
\end{itemize}
Zero-centred outputs matter: a layer whose activations are all positive produces
gradients that share a sign across a whole weight vector, which forces the
optimiser to move in zig-zags.
}

\analogybreak{
No activation is best everywhere. $\tanh$'s stronger central gradient still
saturates; ReLU never saturates on the positive side but discards all negative
information and can die. The right choice depends on where in the network the unit
sits and what range its output must occupy --- which is why the models here use
different activations in different places rather than one throughout.
}

\whyhere{
This is the single highest-impact change made to this project's models. TimeGAN's
embedder, generator and supervisor originally used sigmoid on the latent path. The
diagnostic showed the latent representation stuck near its centre --- mean
$0.5069$, standard deviation $0.04$ --- and the autoencoder reconstruction $R^2$
was approximately zero, meaning Phase 3 was training on noise. Replacing sigmoid
with $\tanh$ on those three latent paths, for exactly the $4\times$ per-layer
reason above, moved the reconstruction $R^2$ to $+0.995$.

The recovery network also uses $\tanh$, but for a different reason: it emits
\emph{data}, and the data was normalised by $\tanh(z/3)$ into $(-1,1)$. Its
activation must match the range of the thing it is reconstructing, which the
round-trip assertion in the code checks to a tolerance of $10^{-3}$.
}

%% ─────────────────────────────────────────────────────────────────────────────
\section{Forward Propagation}
\label{sec:forward}

\situation{
The factory line of \S\ref{sec:composition}: raw material in at one end, finished
product out at the other, each station transforming what it receives.
}

\problem{
With the unit and the non-linearity defined, we need to run a whole network and
keep track of what each stage produced --- because backpropagation will need those
intermediate values again.
}

\workedarith{
A two-layer network: 2 inputs, 2 hidden units with $\tanh$, 1 output, linear.

\[
  W^{(1)} = \begin{bmatrix} 1 & -1 \\ 0.5 & 2 \end{bmatrix} ,
  \quad \mathbf{b}^{(1)} = \begin{bmatrix} 0 \\ 0.5 \end{bmatrix} ,
  \quad W^{(2)} = \begin{bmatrix} 1 & -2 \end{bmatrix} ,
  \quad b^{(2)} = 0.1 .
\]

Input $\mathbf{x} = [1, 2]$.

\medskip
\textbf{Layer 1, pre-activation:}
\[
  \mathbf{z}^{(1)} = W^{(1)}\mathbf{x} + \mathbf{b}^{(1)}
  = \begin{bmatrix} 1(1) + (-1)(2) \\ 0.5(1) + 2(2) \end{bmatrix} + \begin{bmatrix} 0 \\ 0.5 \end{bmatrix}
  = \begin{bmatrix} -1 \\ 5.0 \end{bmatrix} .
\]

\textbf{Layer 1, activation:}
\[
  \mathbf{a}^{(1)} = \tanh\!\begin{bmatrix} -1 \\ 5.0 \end{bmatrix}
  = \begin{bmatrix} -0.7616 \\ 0.9999 \end{bmatrix} .
\]
Note the second unit is saturated: $\tanh(5) = 0.9999$, where
$\tanh' = 1 - 0.9999^2 = 0.0002$. Almost no gradient will pass through it.

\medskip
\textbf{Layer 2:}
\[
  z^{(2)} = 1(-0.7616) + (-2)(0.9999) + 0.1 = -0.7616 - 1.9998 + 0.1 = -2.6614 .
\]

Output $-2.6614$.
}

\formaldef{
For a network of $L$ layers, \textbf{forward propagation} computes
\[
  \mathbf{a}^{(0)} = \mathbf{x} , \qquad
  \mathbf{z}^{(l)} = W^{(l)}\mathbf{a}^{(l-1)} + \mathbf{b}^{(l)} , \qquad
  \mathbf{a}^{(l)} = \sigma^{(l)}\bigl(\mathbf{z}^{(l)}\bigr) ,
\]
for $l = 1,\ldots,L$, with $\mathbf{a}^{(L)}$ the output.

Every $\mathbf{z}^{(l)}$ and $\mathbf{a}^{(l)}$ must be retained: the backward
pass needs them, which is why training uses far more memory than inference.
}

\analogybreak{
The layer-by-layer picture suggests each stage is independently meaningful. It is
not. A trained layer's weights are fitted assuming exactly the distribution the
previous layer sends it, and feeding raw input to a middle layer produces
nonsense. The stations are separable in notation only --- the point already made
in \S\ref{sec:composition}, and the reason a saturated unit like the second one
above corrupts everything downstream of it.
}

\whyhere{
The saturated unit in the worked example is not hypothetical. The diagnostic block
added to TimeGAN's Phase 1 reports exactly this: the fraction of latent
activations with $|h| > 0.98$. Above $50\%$ saturation, $\tanh$ is clipping the
gradient and the layer has effectively stopped learning.
}

%% ─────────────────────────────────────────────────────────────────────────────
\section{Computational Graphs}
\label{sec:compgraph}

\situation{
A recipe written as a flowchart: chop, then fry, then season, then plate. Drawing
the dependencies makes it obvious which steps must precede which, and which could
run in parallel.
}

\problem{
Backpropagation must visit every operation in reverse order and know exactly which
quantities feed which. Doing this by hand for a network with dozens of layers is
error-prone; representing the computation as a graph makes it mechanical.
}

\workedarith{
Take $L = (wx + b - y)^2$ with $w = 2$, $x = 3$, $b = 1$, $y = 5$.

\medskip
\textbf{Forward, as a sequence of elementary operations:}

\medskip
\begin{tabular}{@{}llr@{}}
\toprule
Node & Operation & Value \\
\midrule
$u$ & $w \times x$ & $6$ \\
$v$ & $u + b$      & $7$ \\
$e$ & $v - y$      & $2$ \\
$L$ & $e^2$        & $4$ \\
\bottomrule
\end{tabular}

\medskip
\textbf{Local derivatives, one per edge:}
\[
  \frac{\partial L}{\partial e} = 2e = 4 , \quad
  \frac{\partial e}{\partial v} = 1 , \quad
  \frac{\partial v}{\partial u} = 1 , \quad
  \frac{\partial u}{\partial w} = x = 3 , \quad
  \frac{\partial v}{\partial b} = 1 .
\]

\textbf{Chain them.}
\[
  \frac{\partial L}{\partial w} = 4 \times 1 \times 1 \times 3 = 12 ,
  \qquad
  \frac{\partial L}{\partial b} = 4 \times 1 \times 1 = 4 .
\]

\medskip
\textbf{Verify directly.} $L(w) = (3w + 1 - 5)^2 = (3w-4)^2$, so
$dL/dw = 2(3w-4)(3) = 6(3w-4)$. At $w = 2$: $6(2) = 12$. \checkmark

And $L(b) = (6 + b - 5)^2 = (b+1)^2$, so $dL/db = 2(b+1) = 4$ at $b=1$.
\checkmark

\medskip
Each edge carries one simple derivative; the algorithm just multiplies along
paths and adds across them, exactly as \S\ref{sec:chainmulti} described.
}

\formaldef{
A \textbf{computational graph} is a directed acyclic graph whose nodes are
variables and operations and whose edges are data dependencies. The forward pass
evaluates nodes in topological order.

\textbf{Automatic differentiation} in \emph{reverse mode} traverses the graph
backwards from the output, accumulating $\partial L/\partial(\cdot)$ at each node
by the chain rule. One reverse pass yields the derivative of one output with
respect to \emph{all} inputs, at roughly the cost of the forward pass --- which is
exactly the right trade when there are many parameters and one loss.
}

\analogybreak{
Reverse mode is efficient because there is one output and many inputs. With many
outputs and few inputs, forward mode is cheaper. Neither is universally better;
the asymmetry of the training problem --- millions of parameters, one scalar loss
--- is what makes reverse mode the right choice here.
}

\whyhere{
PyTorch builds this graph as the forward pass executes and \texttt{backward()}
walks it in reverse. Nothing in this project computes a derivative by hand. The
detach operations in the TimeGAN training code exist precisely to cut edges in
this graph, preventing gradients from flowing into components that a given phase
is not supposed to update.
}

%% ─────────────────────────────────────────────────────────────────────────────
\section{Backpropagation}
\label{sec:backprop}

\situation{
A production error appears at the end of the line. To assign responsibility you
walk backwards station by station, asking at each: given the fault that arrived
here, how much of it did this station add?
}

\problem{
The chain rule gives the derivative for one parameter. A network has hundreds of
thousands, and computing them one at a time --- one forward pass per parameter ---
would make training impossible. The whole set must come from a single traversal.
}

\workedarith{
Use the two-layer network of \S\ref{sec:forward} with target $y = -2$. The output
was $z^{(2)} = -2.6614$, so with squared loss
\[
  L = \bigl(z^{(2)} - y\bigr)^2 = (-0.6614)^2 = 0.4374 .
\]

\medskip
\textbf{Step 1, at the output.}
\[
  \frac{\partial L}{\partial z^{(2)}} = 2(-0.6614) = -1.3228 .
\]

\textbf{Step 2, second-layer weights.} Since
$z^{(2)} = W^{(2)}\mathbf{a}^{(1)} + b^{(2)}$,
\[
  \frac{\partial L}{\partial W^{(2)}} = \frac{\partial L}{\partial z^{(2)}}\,\mathbf{a}^{(1)\top}
  = -1.3228\,[-0.7616,\ 0.9999]
  = [1.0075,\ -1.3227] .
\]

\textbf{Step 3, back through $W^{(2)}$.} Multiply by the transpose
(\S\ref{sec:inverse}):
\[
  \frac{\partial L}{\partial \mathbf{a}^{(1)}} = W^{(2)\top}\frac{\partial L}{\partial z^{(2)}}
  = \begin{bmatrix} 1 \\ -2 \end{bmatrix}(-1.3228)
  = \begin{bmatrix} -1.3228 \\ 2.6456 \end{bmatrix} .
\]

\textbf{Step 4, through the activation.} Multiply entrywise by
$\tanh'(z^{(1)}) = 1 - \tanh^2$:
\[
  \tanh'(-1) = 1 - 0.7616^2 = 0.4200 , \qquad
  \tanh'(5) = 1 - 0.9999^2 = 0.0002 .
\]
\[
  \frac{\partial L}{\partial \mathbf{z}^{(1)}}
  = \begin{bmatrix} -1.3228 \times 0.4200 \\ 2.6456 \times 0.0002 \end{bmatrix}
  = \begin{bmatrix} -0.5556 \\ 0.0005 \end{bmatrix} .
\]

\medskip
\textbf{Read the second entry.} The saturated unit receives a gradient of
$0.0005$ --- essentially nothing. Its weights will barely move, however wrong they
are. This is the vanishing gradient, visible in four lines of arithmetic, and it is
the same mechanism that the sigmoid-to-$\tanh$ change in \S\ref{sec:activations}
addressed at a whole-network scale.
}

\formaldef{
\textbf{Backpropagation} computes, for $l = L$ down to $1$:
\[
  \boldsymbol\delta^{(L)} = \nabla_{\mathbf{a}^{(L)}}L \odot \sigma'\bigl(\mathbf{z}^{(L)}\bigr) ,
  \qquad
  \boldsymbol\delta^{(l)} = \Bigl(W^{(l+1)\top}\boldsymbol\delta^{(l+1)}\Bigr) \odot \sigma'\bigl(\mathbf{z}^{(l)}\bigr) ,
\]
where $\odot$ is entrywise multiplication. The parameter gradients are then
\[
  \frac{\partial L}{\partial W^{(l)}} = \boldsymbol\delta^{(l)}\mathbf{a}^{(l-1)\top} ,
  \qquad
  \frac{\partial L}{\partial \mathbf{b}^{(l)}} = \boldsymbol\delta^{(l)} .
\]
Cost: one forward and one backward pass, independent of the parameter count.
}

\analogybreak{
The repeated multiplication by $W^{\top}$ and $\sigma'$ is what makes deep
networks hard to train. If the factors are systematically below 1 the signal
vanishes; above 1 it explodes. Neither is a defect of the chain rule, which is
exact --- both are properties of the arithmetic that implements it, and both are
why initialisation schemes, normalisation layers and gradient clipping exist.
}

\whyhere{
The $\texttt{max\_norm}=1.0$ gradient clip in all three models is the guard against
the exploding case: whenever $\norm{\nabla L}_2$ exceeds 1, the whole gradient is
rescaled. The $\tanh$ change was the guard against the vanishing case. Both are
responses to the same product structure.
}

%% ─────────────────────────────────────────────────────────────────────────────
\section{Training a Network}
\label{sec:training}

\situation{
Practising a piece of music. Play a phrase, hear what went wrong, adjust,
repeat --- thousands of times. No single repetition matters; the accumulation does.
}

\problem{
Forward propagation, backpropagation and gradient descent are three separate
mechanisms. Training is the loop that combines them, and the loop has parameters
of its own that determine whether anything is learned.
}

\workedarith{
\textbf{One update, end to end.} From \S\ref{sec:backprop},
$\partial L/\partial W^{(2)} = [1.0075,\ -1.3227]$ and
$W^{(2)} = [1,\ -2]$. With $\eta = 0.1$:
\[
  W^{(2)} \leftarrow [1,\ -2] - 0.1[1.0075,\ -1.3227] = [0.8993,\ -1.8677] .
\]

\textbf{Did the loss fall?} Recompute the forward pass with the new weights,
holding layer 1 fixed:
\[
  z^{(2)} = 0.8993(-0.7616) + (-1.8677)(0.9999) + 0.1 = -0.6849 - 1.8675 + 0.1 = -2.4524 .
\]
New loss: $(-2.4524 + 2)^2 = 0.2047$, down from $0.4374$. \checkmark More than
halved in one step.

\medskip
\textbf{The loop.}
\begin{enumerate}
  \item Draw a mini-batch of $B$ examples.
  \item Forward pass, storing intermediates.
  \item Compute the loss.
  \item Backward pass for all parameter gradients.
  \item Clip the gradient if $\norm{\nabla L}_2 > 1$.
  \item Update every parameter.
  \item Repeat for the budgeted number of steps.
\end{enumerate}

\medskip
\textbf{Budget, concretely.} This project fixes $9{,}000$ generator updates. With
TimeGAN's 29 batches per epoch (\S\ref{sec:sgd}), that is
$\lceil 9000/29 \rceil = 311$ epochs.
}

\formaldef{
The training loop repeats:
\[
  \theta \leftarrow \theta - \eta\,\mathrm{clip}\!\left(\frac{1}{B}\sum_{i\in\mathcal{B}}\nabla_\theta \ell_i , \; 1.0\right) .
\]
Weights are initialised randomly --- not at zero, which would make every neuron in
a layer compute the same thing and receive the same gradient for ever, leaving them
identical. Standard schemes scale the initial variance by the layer width so that
activations neither shrink nor blow up with depth.
}

\analogybreak{
``Train until convergence'' is not available here. GAN losses oscillate
(\S\ref{sec:minimax}) and a flat loss can signal collapse rather than success, so
there is no reliable stopping signal in the loss itself. Training runs for a fixed
budget and is judged externally.
}

\whyhere{
The fixed-step budget is a methodological requirement, not a convenience. Because
the three models have different batch sizes and different architectures, equal
epoch counts would give unequal gradient steps --- the $29\times$ discrepancy of
\S\ref{sec:sgd}. The parity check in the pipeline asserts equal generator updates
and refuses to run otherwise.
}

%% ─────────────────────────────────────────────────────────────────────────────
\section{Recurrent Networks}
\label{sec:rnn}

\situation{
Reading a sentence. You do not process each word in isolation --- you carry the
meaning built so far and update it with each new word. The carried state is what
lets the ending make sense.
}

\problem{
A feedforward network takes a fixed-size input and has no memory. A return series
is a sequence of arbitrary length in which order is the whole point, and it needs
an architecture that processes it step by step while retaining something.
}

\workedarith{
A minimal recurrent unit: $h_t = \tanh(w_h h_{t-1} + w_x x_t + b)$ with
$w_h = 0.8$, $w_x = 0.5$, $b = 0$, $h_0 = 0$.

Inputs $x = (1, 0, -1)$:
\[
  h_1 = \tanh(0.8(0) + 0.5(1)) = \tanh(0.5) = 0.4621 ,
\]
\[
  h_2 = \tanh(0.8(0.4621) + 0) = \tanh(0.3697) = 0.3537 ,
\]
\[
  h_3 = \tanh(0.8(0.3537) + 0.5(-1)) = \tanh(-0.2170) = -0.2137 .
\]

\medskip
\textbf{Memory is there.} At $t = 2$ the input is $0$, yet $h_2 = 0.3537 \neq 0$:
the state carries information from $t = 1$.

\medskip
\textbf{And it fades.} With no further input, the state would evolve as
$h \leftarrow \tanh(0.8h)$, shrinking each step: $0.3537 \to 0.2755 \to 0.2169$.
The influence of $x_1$ decays geometrically.

\medskip
\textbf{Why that is a problem.} Unrolling over $T$ steps and backpropagating
multiplies $T$ factors of roughly $w_h \tanh'$. With $w_h = 0.8$ and
$\tanh' \leq 1$, over 128 steps --- this project's sequence length --- the product
is at most
\[
  0.8^{128} = 4 \times 10^{-13} .
\]
A dependency spanning the full window transmits essentially no gradient. This is
the vanishing gradient of \S\ref{sec:backprop} in its most acute form.
}

\formaldef{
A \textbf{recurrent neural network} maintains a hidden state:
\[
  \mathbf{h}_t = \sigma\bigl(W_h\mathbf{h}_{t-1} + W_x\mathbf{x}_t + \mathbf{b}\bigr) ,
\]
with the same weights reused at every step --- so the parameter count is
independent of sequence length.

Training uses \textbf{backpropagation through time}: unroll the recurrence into a
$T$-layer feedforward network sharing weights, and backpropagate through it.
}

\analogybreak{
Weight sharing makes RNNs compact and makes them deep in a peculiar way: a
128-step sequence is a 128-layer network in which every layer is identical. The
repeated multiplication by the \emph{same} matrix means the gradient's fate is
governed by that matrix's largest eigenvalue (\S\ref{sec:eigen}) --- below 1 and it
vanishes, above 1 and it explodes. There is no comfortable middle.
}

\whyhere{
TimeGAN is built from recurrent cells. Its embedder, recovery, generator,
supervisor and discriminator are all GRU stacks with hidden dimension 24 and three
layers, operating on sequences of length 128 --- precisely the regime where the
plain RNN above would fail.
}

%% ─────────────────────────────────────────────────────────────────────────────
\section{LSTM and GRU}
\label{sec:gru}

\situation{
Taking notes in a lecture. You do not transcribe everything and you do not
discard everything --- you decide what to write down, what to keep from earlier,
and what is now irrelevant. Selective retention is what makes the notes useful an
hour later.
}

\problem{
The plain RNN forces its state through a multiplication at every step, so
information decays geometrically whether or not it should. What is needed is a
path along which information can travel unchanged, plus a learned decision about
when to change it.
}

\workedarith{
A GRU has an update gate $z_t$ deciding how much of the old state to keep:
\[
  \mathbf{h}_t = (1 - z_t)\odot\mathbf{h}_{t-1} + z_t\odot\tilde{\mathbf{h}}_t .
\]

\medskip
\textbf{Gate closed.} With $z_t = 0$, $\mathbf{h}_t = \mathbf{h}_{t-1}$ exactly.
The state passes through untouched --- no multiplication by a weight matrix, no
decay.

\textbf{Gate open.} With $z_t = 1$, $\mathbf{h}_t = \tilde{\mathbf{h}}_t$: the
old state is discarded entirely.

\textbf{Partly open.} With $z_t = 0.1$, $90\%$ of the old state survives each
step. Over 128 steps that is $0.9^{128} = 1.4\times10^{-6}$ --- still decaying, but
compare the plain RNN's $0.8^{128} = 4\times10^{-13}$. Seven orders of magnitude
better, and $z_t$ is \emph{learned}, so the network can close the gate further
where long memory is needed.

\medskip
\textbf{Why the additive form matters.} The state update is a weighted sum
rather than a matrix product. Differentiating $(1-z)h_{t-1} + z\tilde{h}_t$ with
respect to $h_{t-1}$ gives $(1-z)$ --- a scalar the network controls, not a fixed
eigenvalue it is stuck with.
}

\formaldef{
The \textbf{GRU} (Cho et al.\ 2014) computes
\begin{align*}
  \mathbf{z}_t &= \sigma\bigl(W_z\mathbf{x}_t + U_z\mathbf{h}_{t-1}\bigr) && \text{update gate} \\
  \mathbf{r}_t &= \sigma\bigl(W_r\mathbf{x}_t + U_r\mathbf{h}_{t-1}\bigr) && \text{reset gate} \\
  \tilde{\mathbf{h}}_t &= \tanh\bigl(W\mathbf{x}_t + U(\mathbf{r}_t\odot\mathbf{h}_{t-1})\bigr) && \text{candidate} \\
  \mathbf{h}_t &= (1-\mathbf{z}_t)\odot\mathbf{h}_{t-1} + \mathbf{z}_t\odot\tilde{\mathbf{h}}_t && \text{new state}
\end{align*}
The \textbf{LSTM} (Hochreiter \& Schmidhuber 1997) keeps a separate cell state
with three gates --- input, forget and output --- and roughly a third more
parameters for comparable performance on most tasks.

Note that the gates use $\sigma$ deliberately: a gate must lie in $(0,1)$ to act
as a proportion, which is the one place sigmoid's range is exactly right.
}

\analogybreak{
Gating mitigates vanishing gradients; it does not eliminate them. Very long
dependencies remain hard, and the sequential structure prevents parallelisation
across time --- step $t$ cannot start until $t-1$ finishes. Both limitations are
why convolutional and attention-based architectures displaced recurrent ones for
long sequences.
}

\whyhere{
TimeGAN uses GRU stacks throughout. The activation distinction of
\S\ref{sec:activations} applies precisely here: sigmoid is correct \emph{inside}
the gates, where a proportion in $(0,1)$ is wanted, and wrong on the latent
output path, where the $4\times$ gradient penalty applies with no compensating
benefit. The fix in this project changed the latter and left the former alone.
}

%% ─────────────────────────────────────────────────────────────────────────────
\section{Temporal Convolutions, Dilation and Receptive Fields}
\label{sec:tcn}

\situation{
Skimming a document for a theme. You do not read every word in order; you glance
at every tenth line and still catch the structure. Sampling at intervals covers
far more ground for the same effort.
}

\problem{
Recurrent models are sequential and slow, and their memory decays. Convolutions
process all positions in parallel, but a stack of ordinary convolutions sees only
a narrow window unless it is made impractically deep. Dilation resolves this.
}

\workedarith{
\textbf{Receptive field of a plain stack.} A convolution with kernel size $k = 3$
sees 3 positions. Stacking two extends the view by $(k-1) = 2$ each time, so $L$
layers see $1 + L(k-1)$ positions. To cover 128 days would need
$L = 127/2 \approx 64$ layers.

\medskip
\textbf{With dilation.} A dilated convolution skips $d-1$ positions between taps,
so it spans $(k-1)d + 1$ while using the same $k$ weights. Doubling $d$ each layer
makes the receptive field grow geometrically.

\medskip
\textbf{This project's QuantGAN generator.} Three temporal blocks, each with
\emph{two} causal convolutions of kernel size 3, and dilations $1, 2, 4$:

\medskip
\begin{tabular}{@{}lrrr@{}}
\toprule
Block & Dilation & Extension per block & Cumulative field \\
\midrule
--    & --  & --  & $1$ \\
0     & $1$ & $2 \times 2 \times 1 = 4$  & $5$ \\
1     & $2$ & $2 \times 2 \times 2 = 8$  & $13$ \\
2     & $4$ & $2 \times 2 \times 4 = 16$ & $29$ \\
\bottomrule
\end{tabular}

\medskip
Each block contributes $2 \times (k-1) \times d = 4d$. Total:
\[
  1 + 4(1 + 2 + 4) = 1 + 28 = 29 .
\]
The generator's output at each position depends on the previous \textbf{29} time
steps --- about six trading weeks --- using only three blocks.

\medskip
\textbf{Causality.} A causal convolution pads only on the left, so position $t$
sees $t, t-1, t-2, \ldots$ and never $t+1$. Without this the generator could use
the future to produce the present, which for a time-series model is fatal.
}

\formaldef{
A \textbf{dilated causal convolution} with kernel size $k$ and dilation $d$
computes
\[
  y_t = \sum_{i=0}^{k-1} w_i\,x_{t - i\cdot d} ,
\]
using only indices at or before $t$.

For a network of $L$ blocks with $c$ convolutions each and dilations $d_l$, the
\textbf{receptive field} is
\[
  R = 1 + c(k-1)\sum_{l} d_l .
\]
With $d_l = 2^{l}$ the sum is geometric, so $R$ grows exponentially in depth while
the parameter count grows linearly.
}

\analogybreak{
A large receptive field is an upper bound on what the model \emph{can} see, not
evidence of what it does see. Dilation also leaves gaps: a single dilated layer
samples a sparse subset of positions, and only the stacking of different dilations
fills them in. And a fixed receptive field is a hard horizon --- 29 days here ---
beyond which no dependence can be represented at all, which sits awkwardly against
the long-memory findings of \S\ref{sec:hurst}, where volatility autocorrelation
remains measurable past lag 50.
}

\whyhere{
This is QuantGAN's mechanism. Its generator and discriminator are temporal
convolutional networks with dilations $1, 2, 4$ over channel widths
$[64, 64, 32]$, giving the 29-day receptive field computed above. That is a
concrete, checkable architectural limit, and it is the honest answer to why a
model might reproduce short-horizon volatility clustering while missing the slow
decay that \texttt{hurst\_diff} measures.
}
